\section{Torsion on the full boundary cycle}\label{sec:torsion}

The reduced dlt boundary $\Sigma$ of Corollary~\ref{lem:dlt-boundary}
is a surface to which adjunction and surface abundance apply.  The
prepared-cycle obstruction, however, requires torsion on an effective
Cartier divisor, including its multiplicities.  We first give the
extension argument that passes from the reduced surface to this
possibly nonreduced scheme.

\begin{lemma}[{\cite[Lemma~2.4]{Companion}}]\label{lem:torsion}
Let $T$ be a normal projective threefold over an algebraically closed
field of characteristic $p>0$, and let $E$ be a nonzero effective Cartier
divisor on $T$.  Suppose that $E$ is numerically trivial on every
integral component of its support.  If
\[
 \mathcal O_T(E)|_{E_{\mathrm{red}}\setminus Z}
       \simeq\mathcal O_{E_{\mathrm{red}}\setminus Z}
\]
for a finite set $Z$ of closed points, then $\mathcal O_E(E)$ is torsion.
\end{lemma}

\begin{proof}
We reproduce the proof because both the finite exceptional set and the
nonreduced divisor will occur in our application.  Replace $Z$ by
$Z\cap\operatorname{Supp}E$.  At $z\in Z$, write
$R=\mathcal O_{T,z}$, with maximal ideal $\mathfrak m$, and let $t_z$ be
a local equation of $E$.  Normality gives
\[
 H^1_{\mathfrak m}(R)=0,
 \qquad \ell_R\bigl(H^2_{\mathfrak m}(R)\bigr)<\infty.
\]
Here the second assertion uses the dimension-three hypothesis.
Indeed, present $R=A/I$ with $A$ regular local of dimension $n$ and
essentially of finite type over the field.  Each nonmaximal localization
of $R$ is normal of dimension at most two, hence Cohen--Macaulay.
Catenarity gives
$\dim A_{\mathfrak q}-\dim R_{\mathfrak p}=n-3$ for a prime
$\mathfrak q$ above $\mathfrak p$; the Auslander--Buchsbaum formula
therefore gives projective dimension $n-3$ over $A_{\mathfrak q}$.
Thus
$\operatorname{Ext}^{n-2}_A(R,A)$ is supported only at the closed point.
It is a finite module, so it has finite length; local duality gives the
assertion for $H^2_{\mathfrak m}(R)$.

Choose $j\geq1$ such that $t_z^{j-1}$ annihilates
$H^2_{\mathfrak m}(R)$ at every point of $Z$.  A single $j$ works because
$Z$ is finite.  For every integer $u$, the restriction sequences give
a commutative diagram
\[
\begin{tikzcd}[column sep=small]
0\arrow[r]&\mathcal O_T((u-j)E)\arrow[r]\arrow[d]
 &\mathcal O_T(uE)\arrow[r]\arrow[d,equal]
 &\mathcal O_{jE}(uE)\arrow[r]\arrow[d]&0\\
0\arrow[r]&\mathcal O_T((u-1)E)\arrow[r]
 &\mathcal O_T(uE)\arrow[r]
 &\mathcal O_E(uE)\arrow[r]&0.
\end{tikzcd}
\]
In local trivializations the left vertical map is multiplication by
$t_z^{j-1}$.  Local cohomology with support in the finite set $Z$ is the
direct sum of the corresponding local groups.  The vanishing of the
ambient $H^1_Z$ makes the connecting maps from the $H^1_Z$ of the
rightmost terms to the ambient $H^2_Z$ injective.  Naturality and our
choice of $j$ therefore imply
\begin{equation}\label{eq:torsion-obstruction-zero}
 H^1_Z\bigl(T,\mathcal O_{jE}(uE)\bigr)
    \longrightarrow H^1_Z\bigl(T,\mathcal O_E(uE)\bigr)
       \quad\text{is zero for every }u\in\mathbb Z.
\end{equation}
The same $j$ works for every twist: locally the invertible sheaves in
the diagram are free of rank one, and the annihilating map is still
multiplication by $t_z^{j-1}$.

Now fix this thickening $jE$.  The ideal of
$E_{\mathrm{red}}\setminus Z$ in $jE\setminus Z$ is nilpotent.  Lift a
chosen generator of $\mathcal O_T(E)|_{E_{\mathrm{red}}\setminus Z}$ to
local generating sections on a finite open cover of $jE\setminus Z$.
On overlaps the lifts differ by sections in that nilpotent ideal times
the line bundle.  For $u=p^e$ sufficiently large, their $u$th tensor
powers agree: in local frames this is the identity
$(a+b)^{p^e}=a^{p^e}+b^{p^e}$, with $b^{p^e}=0$.  They consequently
glue to a generator
\[
 s_j\in H^0\bigl(jE\setminus Z,\mathcal O_{jE}(uE)\bigr).
\]
Let $s$ be its restriction to $E\setminus Z$.  The obstruction to
extending $s$ over $E$ is the image of the obstruction to extending
$s_j$ over $jE$.  Equation~\eqref{eq:torsion-obstruction-zero}, applied
to the exact sequences for sections with support in $Z$, kills that
image.  Hence $s$ extends to a section of $\mathcal O_E(uE)$.

This extended section has no zeros.  On each integral projective
surface component $F$ of $E_{\mathrm{red}}$, it is nonzero because it is
a generator off $Z$.  If its restriction to $F$ vanished anywhere,
its zero scheme would be a nonempty effective Cartier divisor on $F$,
and would have positive intersection with an ample divisor.  This
contradicts the numerical triviality of $\mathcal O_F(uE)$.  The
argument uses only integrality of $F$, not normality.  The section is
therefore a generator on the reduction, and then on $E$, since an
element invertible modulo a nilpotent ideal is invertible.  Thus
$\mathcal O_E(uE)\simeq\mathcal O_E$.
\end{proof}

We apply this lemma to the reduced dlt model obtained by isolation.
Recall that $s$ is the multiplicity of $S_Y$ in $N_Y$ and $\tau>0$ is
the rational number satisfying $K_Y+G_Y\sim_{\mathbb Q}\tau N_Y$.

\begin{proposition}\label{prop:boundary-torsion}
For the projective crepant dlt modification
$f:(Y_d,\Sigma)\to(Y,S_Y)$ of Corollary~\ref{lem:dlt-boundary}, there is a
positive Cartier multiple $M$ of $f^*S_Y$ such that
\[
 \operatorname{Supp}M=\Sigma,
 \qquad M|_{\Sigma_i}\equiv0\quad\text{for every component }\Sigma_i,
 \qquad \mathcal O_M(M)\text{ is torsion}.
\]
\end{proposition}

\begin{proof}
The relations furnished by isolation are
\begin{equation}\label{eq:adjunction-input}
 K_{Y_d}+\Sigma\sim_{\mathbb Q}\tau s f^*S_Y
       \quad\text{on a neighborhood of }\Sigma,
 \qquad f^*S_Y|_{\Sigma_i}\equiv0.
\end{equation}
Let $\sigma:\Sigma'\to\Sigma$ be the finite $S_2$-ification.  It is
unchanged away from a finite set $Z\subset\Sigma$
\cite[Proposition~2.14]{XuNuTwo}.  Reduced-boundary adjunction gives
an effective rational divisor $\Delta'$ such that $(\Sigma',\Delta')$
is a projective semi-log-canonical surface pair and
\begin{equation}\label{eq:slc-adjunction}
 K_{\Sigma'}+\Delta'
       \sim_{\mathbb Q}\sigma^*((K_{Y_d}+\Sigma)|_\Sigma)
       \sim_{\mathbb Q}\tau s\,\sigma^*(f^*S_Y|_\Sigma).
\end{equation}
We use adjunction in the form of
\cite[Lemma~2.19]{XuNuTwo}; see also
\cite[Section~2.5]{Companion}.  The $S_2$-ification is needed because
the reduced boundary itself need not satisfy $S_2$.  Its
codimension-one singularities are nodes, and on its normalization the
adjunction boundary includes the conductor with coefficient one;
the normalized pair is log canonical.  These are precisely the
semi-log-canonical conditions needed for surface abundance.

The surface $\Sigma'$ is projective because it is finite over $\Sigma$.
The last class in \eqref{eq:slc-adjunction} is numerically trivial,
by \eqref{eq:adjunction-input} and the projection formula for curves.
Posva's abundance theorem for projective semi-log-canonical surfaces
in positive characteristic therefore makes
$K_{\Sigma'}+\Delta'$ semiample
\cite[Theorem~3.1]{Posva}.  A semiample numerically trivial rational
line bundle is $\mathbb Q$-linearly trivial: a globally generated Cartier
multiple defines a morphism contracting every curve, hence constant
on each connected component.  Its line bundle is trivial there.

Since $\tau s$ is positive and rational, clearing denominators in
\eqref{eq:slc-adjunction} yields a positive Cartier multiple $M$ of
$f^*S_Y$ with $\sigma^*\mathcal O_\Sigma(M)$ trivial.  In particular,
\[
 \mathcal O_{Y_d}(M)|_{\Sigma\setminus Z}
        \simeq\mathcal O_{\Sigma\setminus Z}.
\]
The divisor $M$ is effective, its reduction is $\Sigma$, and it is
numerically trivial on every component of its support.  Applying
Lemma~\ref{lem:torsion} on the normal projective threefold $Y_d$ proves
torsion on the full Cartier scheme $M$.
\end{proof}
